Trong phần này bạn tạo nơi lưu và nơi phân phối giao diện web. Bucket S3
giữ file, CloudFront đưa file đó tới người dùng.

Điểm cần chú ý: bucket sẽ được để \textbf{riêng tư hoàn toàn}. Người
dùng không truy cập thẳng vào S3 mà chỉ đi qua CloudFront. Cơ chế cho
phép điều này gọi là \textbf{Origin Access Control}.

\vspace{0.5em}\noindent\textbf{Bước 1. Tạo bucket lưu giao
diện}\par

\begin{enumerate}
\def\labelenumi{\arabic{enumi}.}
\tightlist
\item
  Mở
  \href{https://s3.console.aws.amazon.com/s3/home?region=ap-southeast-1}{Amazon
  S3 console}
\item
  Chọn \textbf{Create bucket}
\item
  Điền thông tin:

  \begin{itemize}
  \tightlist
  \item
    \textbf{Bucket name}:
    \texttt{fcj-recsys-frontend-\textless{}tên-bạn\textgreater{}} (tên
    bucket là duy nhất toàn cầu, thêm hậu tố để tránh trùng)
  \item
    \textbf{AWS Region}: Asia Pacific (Singapore)
    \texttt{ap-southeast-1}
  \item
    \textbf{Block Public Access}: \textbf{giữ nguyên bật tất cả}
  \end{itemize}
\item
  Chọn \textbf{Create bucket}
\end{enumerate}

\begin{figure}
\centering
\pandocbounded{\includegraphics[keepaspectratio,alt={Tạo bucket}]{5-Workshop/5.3/create-bucket.png}}
\caption{Tạo bucket}
\end{figure}

Nhiều hướng dẫn cũ bảo bật Static website hosting và mở public cho
bucket. Cách đó không còn được khuyến nghị vì lộ bucket ra Internet. Ở
workshop này ta giữ bucket kín và để CloudFront làm cửa duy nhất.

\vspace{0.5em}\noindent\textbf{Bước 2. Build giao
diện}\par

Tại máy của bạn, vào thư mục \texttt{frontend} và build:

\begin{Shaded}
\begin{Highlighting}[]
\BuiltInTok{cd}\NormalTok{ frontend}
\ExtensionTok{npm}\NormalTok{ install}
\ExtensionTok{npm}\NormalTok{ run build}
\end{Highlighting}
\end{Shaded}

Lệnh này tạo ra thư mục \texttt{dist} chứa toàn bộ file tĩnh.

Chưa cần lo địa chỉ API lúc này. Sau khi tạo xong API Gateway ở mục 5.5,
bạn sẽ quay lại đặt biến \texttt{VITE\_API\_BASE\_URL} rồi build lại.

\vspace{0.5em}\noindent\textbf{Bước 3. Tải file lên
S3}\par

\begin{Shaded}
\begin{Highlighting}[]
\ExtensionTok{aws}\NormalTok{ s3 sync dist/ s3://fcj{-}recsys{-}frontend{-}}\OperatorTok{\textless{}}\NormalTok{tên{-}bạn}\OperatorTok{\textgreater{}}\NormalTok{/ }\AttributeTok{{-}{-}delete}
\end{Highlighting}
\end{Shaded}

Tham số \texttt{-\/-delete} xoá những file cũ không còn trong bản build
mới, tránh để lại rác qua nhiều lần deploy.

Kiểm tra lại trên console, bucket phải có \texttt{index.html} cùng thư
mục \texttt{assets}.

\vspace{0.5em}\noindent\textbf{Bước 4. Tạo CloudFront
distribution}\par

\begin{enumerate}
\def\labelenumi{\arabic{enumi}.}
\tightlist
\item
  Mở \href{https://console.aws.amazon.com/cloudfront/v4/home}{CloudFront
  console}
\item
  Chọn \textbf{Create distribution}
\item
  Ở phần \textbf{Origin}:

  \begin{itemize}
  \tightlist
  \item
    \textbf{Origin domain}: chọn bucket vừa tạo từ danh sách
  \item
    \textbf{Origin access}: chọn \textbf{Origin access control settings
    (recommended)}
  \item
    Chọn \textbf{Create new OAC}, giữ nguyên tuỳ chọn mặc định, bấm
    \textbf{Create}
  \end{itemize}
\item
  Ở phần \textbf{Default cache behavior}:

  \begin{itemize}
  \tightlist
  \item
    \textbf{Viewer protocol policy}: \textbf{Redirect HTTP to HTTPS}
  \item
    \textbf{Allowed HTTP methods}: \texttt{GET,\ HEAD}
  \end{itemize}
\item
  Ở phần \textbf{Settings}:

  \begin{itemize}
  \tightlist
  \item
    \textbf{Default root object}: \texttt{index.html}
  \end{itemize}
\item
  Chọn \textbf{Create distribution}
\end{enumerate}

\begin{figure}
\centering
\pandocbounded{\includegraphics[keepaspectratio,alt={Tạo distribution}]{5-Workshop/5.3/create-distribution.png}}
\caption{Tạo distribution}
\end{figure}

\textbf{Bảo mật nâng cao:} Ở bước tạo CloudFront này, bạn có thể áp dụng
thêm \textbf{AWS WAF} (Web Application Firewall) để chặn Botnet và các
cuộc tấn công DDoS vào giao diện web của bạn.

\vspace{0.5em}\noindent\textbf{Bước 5. Cấp quyền cho CloudFront đọc
bucket}\par

Sau khi tạo xong, CloudFront hiện một banner nhắc cập nhật bucket
policy. Chọn \textbf{Copy policy}, rồi:

\begin{enumerate}
\def\labelenumi{\arabic{enumi}.}
\tightlist
\item
  Sang
  \href{https://s3.console.aws.amazon.com/s3/home?region=ap-southeast-1}{S3
  console}, mở bucket
\item
  Vào tab \textbf{Permissions}
\item
  Ở mục \textbf{Bucket policy}, chọn \textbf{Edit}, dán policy vừa copy,
  chọn \textbf{Save changes}
\end{enumerate}

Policy này chỉ cho phép đúng distribution vừa tạo đọc bucket, không mở
cho ai khác.

\vspace{0.5em}\noindent\textbf{Bước 6. Xử lý lỗi 403 với ứng dụng một
trang}\par

Giao diện React dùng định tuyến phía client. Khi người dùng mở thẳng
\texttt{/cart}, CloudFront đi tìm file tên \texttt{cart} trong S3, không
thấy nên trả về lỗi. Cần chuyển các lỗi đó về \texttt{index.html}.

\begin{enumerate}
\def\labelenumi{\arabic{enumi}.}
\tightlist
\item
  Trong distribution, mở tab \textbf{Error pages}
\item
  Chọn \textbf{Create custom error response}, khai báo:

  \begin{itemize}
  \tightlist
  \item
    \textbf{HTTP error code}: \texttt{403\ Forbidden}
  \item
    \textbf{Customize error response}: Yes
  \item
    \textbf{Response page path}: \texttt{/index.html}
  \item
    \textbf{HTTP Response code}: \texttt{200\ OK}
  \end{itemize}
\item
  Lặp lại một lần nữa cho mã lỗi \texttt{404\ Not\ Found}
\end{enumerate}

\vspace{0.5em}\noindent\textbf{Kiểm tra}\par

Chờ khoảng 5 phút cho distribution chuyển sang trạng thái
\textbf{Deployed}, sau đó mở địa chỉ ở cột \textbf{Distribution domain
name}, dạng \texttt{dxxxxxxxxx.cloudfront.net}.

Giao diện phải hiển thị được. Phần dữ liệu sản phẩm chưa có vì API chưa
tạo, điều đó bình thường.

\textbf{Ghi nhớ điều này, nó sẽ tiết kiệm cho bạn hàng giờ.} Mỗi lần tải
bản build mới lên S3, bạn \textbf{bắt buộc} phải tạo invalidation, nếu
không CloudFront vẫn phục vụ bản cũ trong bộ nhớ đệm và bạn sẽ tưởng mã
nguồn bị lỗi.

\begin{Shaded}
\begin{Highlighting}[]
\ExtensionTok{aws}\NormalTok{ cloudfront create{-}invalidation }\AttributeTok{{-}{-}distribution{-}id} \OperatorTok{\textless{}}\NormalTok{ID{-}CUA{-}BAN}\OperatorTok{\textgreater{}}\NormalTok{ {-}{-}paths }\StringTok{"/*"}
\end{Highlighting}
\end{Shaded}
